% !TeX root = ../main.tex
\section{Цель работы и исходные данные}

Целью данной практической работы является исследование устройства цифровой квадратурной демодуляции радиосигналов с промежуточной частоты на видеочастоту, а также анализ влияния разрядности гетеродина и амплитудно-фазового разбаланса квадратурных каналов на спектральные характеристики выходного сигнала. В среде MATLAB выполняется имитационное моделирование тракта дискретизации и цифрового гетеродинирования с последующей оценкой побочных спектральных составляющих.

В соответствии с заданием, параметры устройства рассчитываются исходя из номера по журналу ($N_{\text{пж}} = 8$) с учётом отстройки частоты входного монохроматического сигнала на $0{,}1$~МГц (Таблица~\ref{tab:initial_data}).

\begin{table}[H]
    \small
    \caption{Параметры устройства цифровой демодуляции (вариант 8)}
    \label{tab:initial_data}
    \begin{tabularx}{\textwidth}{|X|c|c|}
        \hline
        \textbf{Наименование параметра} & \textbf{Формула} & \textbf{Значение} \\ \hline
        Промежуточная частота входного сигнала, $f_C$ & $130 + N_{\text{пж}} \cdot 10 + 0{,}1$ & $210{,}1$ МГц \\ \hline
        Частота дискретизации АЦП, $f_S$ & $100 + N_{\text{пж}} \cdot 10$ & $180{,}0$ МГц \\ \hline
        Количество бит АЦП, $n_a$ & $5 + N_{\text{пж}}$ & 13 бит \\ \hline
        Центральная частота цифрового гетеродина (ЦГ), $f_G$ & $130 + N_{\text{пж}} \cdot 10$ & $210{,}0$ МГц \\ \hline
        Количество бит цифрового гетеродина, $n_g$ & $5 + N_{\text{пж}}$ & 13 бит \\ \hline
    \end{tabularx}
\end{table}
